\documentclass[11pt]{amsart}
\usepackage[T1]{fontenc}
\usepackage[utf8]{inputenc}
\usepackage{amssymb}
\usepackage[margin=1.05in]{geometry}
\usepackage[hidelinks]{hyperref}
\newcommand{\stamptime}{\the\numexpr\time/60\relax:\ifnum\numexpr\time-60*(\time/60)\relax<10 0\fi\the\numexpr\time-60*(\time/60)\relax}
\AddToHook{shipout/foreground}{\put(0.5\paperwidth,-\paperheight+0.55cm){\makebox[0pt]{\tiny\ttfamily unit/doc/\jobname.tex \ --- \ compiled \today, \stamptime}}}
\newtheorem{theorem}{Theorem}
\newtheorem{proposition}[theorem]{Proposition}
\newtheorem{lemma}[theorem]{Lemma}
\newtheorem{corollary}[theorem]{Corollary}
\theoremstyle{definition}
\newtheorem{definition}[theorem]{Definition}
\theoremstyle{remark}
\newtheorem{remark}[theorem]{Remark}
\newcommand{\Om}{\Omega}
\newcommand{\Map}{\mathrm{Map}_*}
\newcommand{\fib}{\mathrm{fib}}
\newcommand{\Z}{\mathbb{Z}}
\newcommand{\N}{\mathbb{N}}
\newcommand{\code}[1]{\texttt{#1}}
\newcommand{\cT}{\mathsf{isTrunc}}
\newcommand{\cC}{\mathsf{isConn}}
\newcommand{\band}{\mathsf{band}}
\newcommand{\Lp}{\mathsf{L}}
\newcommand{\cc}{\cdot}
\newcommand{\refl}{\mathsf{refl}}
\newcommand{\ev}{\mathsf{ev}}
\title{The circle property of $K(\Z,1)$, and the vanishing of $[K(\Z,2),K(B,3)]_*$}
\date{28 October 2026}
\begin{document}
\maketitle

\noindent\textbf{Purpose and status.} Informal proof for the next step of
\code{sag/pi3-s2-postnikov} (handoff \code{doc/pi3-s2-postnikov-handoff-2.txt},
\S3): maps out of $K(\Z,1)$ are determined by a point and a loop, and
the obligation (H3) of \code{KZ2Cohomology} follows. Written before the
formal proof; each numbered statement is a declaration
(\S\ref{sec:corr}). It was refereed independently before formalisation: no
mathematical error was found; the corrections asked for (directions of
retracts, lemmas missing from the library, the step on paths in $\Lp Y$)
are included below. \emph{Status: checked.} \code{CircleProperty} and
\code{H3Proof} (63 modules, 0 failed); \code{H3Proof.h3 : KZ2Cohomology.H3}.
Every equivalence used below is either given or obtained as a retract; all
transfers along an equivalence are proved by induction on it
(\code{Torsor.eqind}), so no half-adjoint coherence is needed. Foundations:
those of the development (univalence, images and their surjectivity;
function extensionality derived); no higher inductive type and no new
axiom. The argument follows the idea of Bezem--Buchholtz--Grayson--Shulman
(\emph{Construction of the circle in UniMath}, arXiv:1910.01856): a map out
of the torsors is defined by contracting, at every point, a type of
equivariant data. Here the torsor structure is used only through the
band, so the statement is about an abstract pointed connected type with
loop group $\Z$.

\section{Setting and conventions}

Paths compose in diagrammatic order: for $p:x=y$, $q:y=z$,
$p\cc q:x=z$ is \code{trans p q}, defined by $J$ on $q$; hence
$p\cc\mathsf{refl}\equiv p$ and $\mathsf{ap}_F\,\mathsf{refl}\equiv\mathsf{refl}$
hold by definition, while $\mathsf{refl}\cc p=p$ is a lemma.
$\Z$ is \code{Integers.Int} with $\mathsf{suc}$, $+$ (recursive in the
second argument, so $n+1\equiv\mathsf{suc}\,n$ by definition).

$K$ is $K(\Z,1)$ of the library: $K=\code{EM1[Zab].K1}$, the small
$\Z$-torsors, pointed by $k_0$ (the trivial torsor). We use only:
\begin{itemize}
\item[(K1)] $K$ is connected: $\cC\,2\,K$ (\code{conn1}); so a family of
  propositions over $K$ has a section as soon as it holds at $k_0$
  (\code{Connected.conn-elim}, $n=0$);
\item[(K2)] the band $\band:(k_0=k_0)\to\Z$ is an equivalence
  (\code{band-equiv}) and a homomorphism,
  $\band(p\cc q)=\band\,p+\band\,q$ (\code{band-hom}).
\end{itemize}
Let $g:k_0=k_0$ be the generator, $g=\band^{-1}(1)$, so $\band\,g=1$.
Then, by (K2) and commutativity,
\[(\ast)\qquad \band(g\cc p)=1+\band\,p=\band\,p+1\equiv\mathsf{suc}(\band\,p).\]
For a type $Y$ the free loop type is $\Lp Y=\Sigma(y:Y).\,y=y$, and
\[E:(K\to Y)\to\Lp Y,\qquad E\,F=(F\,k_0,\ \mathsf{ap}_F\,g).\]

\section{Recursion on $\N$ and on $\Z$}

\begin{lemma}[recursion on $\N$]\label{lem:N}
Let $W$ be a type, $\psi:W\to W$ any map and $w:W$. The type
\[C_0(\psi,w)=\Sigma(u:\N\to W).\,(u\,0=w)\times\Pi(n:\N).\,u(n+1)=\psi(u\,n)\]
is contractible.
\end{lemma}
\begin{proof}
Centre: $(r,\mathsf{refl},\lambda n.\mathsf{refl})$ with $r$ defined by
recursion, $r\,0=w$, $r(n+1)=\psi(r\,n)$. Given $(u,e_0,e)$, define
$H:\Pi n.\,r\,n=u\,n$ by recursion: $H\,0=e_0^{-1}$,
$H(n+1)=\mathsf{ap}_\psi(H\,n)\cc e(n)^{-1}$. A path
$(r,\mathsf{refl},\lambda n.\mathsf{refl})=(u,e_0,e)$ is given by
$\mathsf{funext}\,H$ and the two transports, computed by $J$ on a general
path $P:r=u$ (then specialised with $\mathsf{happly}(\mathsf{funext}\,H)=H$):
the transport of $\mathsf{refl}:r\,0=w$ is $(\mathsf{happly}\,P\,0)^{-1}$,
equal to $e_0$ for $P=\mathsf{funext}\,H$; the transport of
$\lambda n.\mathsf{refl}$ at $n$ is
$(\mathsf{happly}\,P\,(n+1))^{-1}\cc\mathsf{ap}_\psi(\mathsf{happly}\,P\,n)$,
which for $P=\mathsf{funext}\,H$ is
$(\mathsf{ap}_\psi(H\,n)\cc e(n)^{-1})^{-1}\cc\mathsf{ap}_\psi(H\,n)=e(n)$.
\end{proof}

\begin{lemma}[recursion on $\Z$]\label{lem:Z}
Let $\varphi:W\to W$ be an equivalence and $w:W$. The type
\[C_1(\varphi,w)=\Sigma(h:\Z\to W).\,(h\,0=w)\times
  \Pi(n:\Z).\,h(\mathsf{suc}\,n)=\varphi(h\,n)\]
is contractible.
\end{lemma}
\begin{proof}
Write $h^+(n)=h(\mathsf{pos}\,n)$, $h^-(m)=h(\mathsf{negsuc}\,m)$. A
function on $\Z$, and a $\Pi$ over $\Z$, is the pair of its restrictions
(by cases on $\Z$, and funext). Since $\mathsf{suc}(\mathsf{pos}\,n)\equiv
\mathsf{pos}(n+1)$, $\mathsf{suc}(\mathsf{negsuc}\,0)\equiv\mathsf{pos}\,0$
and $\mathsf{suc}(\mathsf{negsuc}(m+1))\equiv\mathsf{negsuc}\,m$, the
condition on $h$ is the conjunction of
$h^+(n+1)=\varphi(h^+n)$, $h^+0=\varphi(h^-0)$ and
$h^-m=\varphi(h^-(m+1))$. Hence
\[C_1(\varphi,w)\simeq\Sigma(h^+,e_0,e^+):C_0(\varphi,w).\ B(h^+0),\qquad
B(v)=\Sigma(h^-:\N\to W).\,(v=\varphi(h^-0))\times
  \Pi m.\,h^-m=\varphi(h^-(m+1)).\]
For an equivalence $\varphi$ with inverse $\varphi^{-1}$ (its section),
$(a=\varphi\,b)$ is a retract of $(b=\varphi^{-1}a)$: by induction on
$\varphi$ it suffices to take $\varphi=\mathrm{id}$, where it is inversion
(\code{equiv-rt}). So $B(v)$ is a retract of $C_0(\varphi^{-1},\varphi^{-1}v)$,
contractible by Lemma~\ref{lem:N}; and $C_0(\varphi,w)$ is contractible by
Lemma~\ref{lem:N}. A $\Sigma$ of contractible types over a contractible
type is contractible.

\emph{As formalised} (the ``$\simeq$'' above taken as retracts in the right
direction): $C_1$ is fibrewise a retract of $\Sigma h.\,\mathrm{Data}'(h)$,
the four conditions on $h$ by cases; this is $\Sigma h.\,\Phi(\mathrm{res}\,h)$
for the restriction $\mathrm{res}:(\Z\to W)\simeq(\N\to W)^2$, hence a retract
of $\Sigma x.\,\Phi\,x$ (\code{Retracts.rt-base}, $\Sigma$ over the domain of
the equivalence), which is $\Sigma(a:C_0(\varphi,w)).\,B(a)$ by reassociation.
\end{proof}

\section{The circle property}

\begin{definition}
For $(y,l):\Lp Y$ and $k:K$ let
\[Q_k=\Sigma(y':Y).\,\Sigma\big(h:(k_0=k)\to(y=y')\big).\,
  \Pi(p:k_0=k).\ h(g\cc p)=l\cc h(p).\]
\end{definition}
\noindent(The type $k_0=k$ is the $\Z$-torsor at $k$, acted on by
$g\cc-$; $Q_k$ is the type of points $y'$ with an equivariant map from it
to the paths $y=y'$, $\Z$ acting by $l\cc-$.)

\begin{lemma}\label{lem:Q}
$Q_k$ is contractible for every $k:K$.
\end{lemma}
\begin{proof}
Contractibility is a proposition, so by (K1) it suffices to treat
$k=k_0$. Let $P:((k_0=k_0)\to W)\simeq(\Z\to W)$ be precomposition with
$\band^{-1}$, whose inverse is $\tilde h\mapsto\tilde h\circ\band$. A
$\Sigma$ over the domain of an equivalence is a retract of the $\Sigma$ over
its codomain of the family composed with the inverse (\code{rt-base'}, by
induction on $P$); so $\Sigma h.\,\mathrm{Cond}(h)$ is a retract of
$\Sigma\tilde h.\,\mathrm{Cond}(\tilde h\circ\band)$. For $h=\tilde h\circ\band$
the condition reads $\Pi p.\,\tilde h(\band(g\cc p))=l\cc\tilde h(\band\,p)$,
a retract (along $(\ast)$, pointwise) of
$\Pi p.\,\tilde h(\mathsf{suc}(\band\,p))=l\cc\tilde h(\band\,p)$, itself a
retract of $\Pi n.\,\tilde h(\mathsf{suc}\,n)=l\cc\tilde h(n)$ ($\Pi$ over the
equivalence $\band$, \code{pi-rt}, by induction on it). Hence, adding the
singleton $q=\tilde h(0)$, $Q_{k_0}$ is a retract of
\[\Sigma(y':Y).\,\Sigma(q:y=y').\ C_1(l\cc-,\,q)\]
with $W=(y=y')$ and $\varphi=l\cc-$, an equivalence. The base
$\Sigma(y':Y).\,y=y'$ is contractible and every $C_1(l\cc-,q)$ is
contractible (Lemma~\ref{lem:Z}).
\end{proof}

\begin{theorem}[circle property]\label{thm:circle}
For every type $Y$, $E:(K\to Y)\to\Lp Y$ is an equivalence.
\end{theorem}
\begin{proof}
We show that every fibre $\fib_E(y,l)$ is a retract of
$\Pi(k:K).\,Q_k$, which is contractible by Lemma~\ref{lem:Q} and
function extensionality.
\begin{enumerate}
\item Paths in $\Lp Y$: $\big((a,u)=(b,v)\big)$ is a retract of
  $\Sigma(e:b=a).\ e\cc u=v\cc e$ (\code{lp-rt}): one way by $J$ on the
  path in $\Lp Y$, giving $(\mathsf{refl},\mathsf{refl}\cc u=u)$ at
  $\mathsf{refl}$; the other by based $J$ on $e$, giving at $e=\mathsf{refl}$ the
  path $(a,u)=(a,v)$ from $u=\mathsf{refl}\cc u=v$; the composite is the
  identity by $J$ and $\mathsf{trans}(p^{-1},p)=\mathsf{refl}$. So
  \[\fib_E(y,l)\ \triangleleft\ \Sigma(F:K\to Y).\,\Sigma(e:y=F\,k_0).\ e\cc\mathsf{ap}_F\,g=l\cc e.\]
\item Distributivity of $\Pi$ over $\Sigma$ (an equivalence whose two
  composites are the identity by definition):
  \[\Pi_k\,Q_k\simeq\Sigma(F:K\to Y).\,\Sigma\big(H:\Pi_k\,(k_0=k)\to(y=F\,k)\big).\,
    \Pi_k\Pi_p\ H_k(g\cc p)=l\cc H_k(p).\]
\item Yoneda: $j:(y=F\,k_0)\to\Pi_k\,(k_0=k)\to(y=F\,k)$,
  $j(e)_k(p)=e\cc\mathsf{ap}_F\,p$, is an equivalence: its inverse is
  evaluation at $(k_0,\mathsf{refl})$, with $j(e)_{k_0}(\mathsf{refl})\equiv e$,
  and $j(H_{k_0}\mathsf{refl})_k(p)=H_k(p)$ by $J$ on $p$. A $\Sigma$ over a
  base $X$ is a retract of the $\Sigma$ over a base equivalent to $X$
  (\code{Retracts.rt-base}), so the type of (2) has as a retract
  \[\Sigma F.\,\Sigma(e:y=F\,k_0).\,\Pi_k\Pi_p\
    e\cc\mathsf{ap}_F(g\cc p)=l\cc(e\cc\mathsf{ap}_F\,p).\]
\item Based path induction: for a family $D$ over $\Sigma(k:K).\,k_0=k$,
  $D(k_0,\mathsf{refl})$ is a retract of $\Pi_k\Pi_p\,D(k,p)$ (extension by
  $J$, evaluation at $\mathsf{refl}$; the composite is the identity by
  definition). This is the general lemma \code{yon-rt}, of which
  \code{DeloopMaps.Maps.Over.Fixed.yoneda} is an instance.
  At $(k_0,\mathsf{refl})$ the condition of (3) is, by definition,
  $e\cc\mathsf{ap}_F\,g=l\cc e$.
\end{enumerate}
Composing, $\fib_E(y,l)$ (by (1)) is a retract of the type of (3), hence
of $\Pi_k\,Q_k$, which is contractible; so $E$ has contractible fibres
(\code{Whitehead.contr-equiv}).
\end{proof}

\begin{corollary}[pointed form]\label{cor:pointed}
For a pointed type $(Y,y_0)$ the canonical pointed evaluation
$\ev_\circ:\Map(K,Y)\to(y_0=y_0)$, $F\mapsto(\mathsf{loop}\,F)\,g$, is an
equivalence, and $\ev_\circ(F,p_F)=p_F^{-1}\cc\mathsf{ap}_F\,g\cc p_F$
(\code{DeloopMaps.Maps.loop-ap}).
\end{corollary}
\begin{proof}
$\Map(K,Y)=\Sigma F.\,F\,k_0=y_0$ is the fibre over $y_0$ of
$\pi\circ E$, with $\pi:\Lp Y\to Y$ the first projection, since
$\pi(E\,F)\equiv F\,k_0$. As $E$ is an equivalence, $(F,p)\mapsto(E\,F,p)$
is an equivalence onto $\fib_\pi(y_0)$ (by induction on $E$); and
$((y,l),p)\mapsto p_*l$ is an equivalence $\fib_\pi(y_0)\simeq(y_0=y_0)$,
with inverse $l\mapsto((y_0,l),\refl)$ (both round trips by $J$). The
composite sends $(F,p_F)$ to $(p_F)_*(\mathsf{ap}_F\,g)$, and
$p_*u=p^{-1}\cc u\cc p$ for a loop $u$ (by $J$ on $p$), which is the
formula; so $\ev_\circ$ is homotopic to an equivalence. \emph{Update
(\S2a of handoff 3):} first formalised as a retract only; now the
equivalence (\code{Circle.pointed-equiv}). A retract of a connected type
is connected (\code{conn-rt}).
\end{proof}
\noindent For (H3) only a consequence is used: if $(y_0=y_0)$ is
connected, so is $\Map(K,Y)$ (\code{Deloop.equiv-conn}).

\section{(H3)}

\begin{theorem}\label{thm:H3}
For every abelian group $B$, $[K(\Z,2),K(B,3)]_*$ is contractible
(\code{KZ2Cohomology.H3}).
\end{theorem}
\begin{proof}
Write $A_2=K(\Z,2)$, $B_3=K(B,3)$ (\code{KK[Zab].K two},
\code{KK[B].K three}), with base points $a_0$, $b_0$, and $0=(\lambda
x.\,b_0,\mathsf{refl})$ the constant pointed map. The set of classes is
inhabited by $[0]$; by induction on classes (into the propositions $[f]=[0]$)
it suffices that every $f:A_2\to_*B_3$ is merely equal to $0$.

\emph{The looping has propositional fibres.} For
$\lambda=\code{Maps[A2,a0,B3,b0].loop}:\Map(A_2,B_3)\to\Map(\Om A_2,\Om B_3)$ the
fibre over $\phi$ is $\code{Unloop}\,\phi=\Sigma F.\,\phi=\lambda F$.
By Corollary~10 of \code{DeloopMaps} (\code{Levels[c,t].cor10}) with
$c=1$, $t=1$: $A_2$ is $\cC\,3$ (\code{KK.conn 1}) and $B_3$ is
$\cT\,5=\cT(2+(2+1))$ (\code{KK.trunc 2}); so $\code{Unloop}\,\phi$ is
$\cT\,1$, hence a proposition (\code{HLevels.trunc-one-prop}). (Orientation:
$\code{Unloop}\,\phi=\Sigma F.\,\phi=\lambda F$, not $\fib_\lambda$; only
\code{Unloop} is used.)

\emph{The target of the looping is connected.} The pointed equivalences
$e_1:\Om A_2\simeq_*K(\Z,1)$ and $e_2:\Om B_3\simeq_*K(B,2)$ (\code{KK.e}) give,
by pre- and post-composition, $\Map(\Om A_2,\Om B_3)$ as a retract of
$\Map(K(\Z,1),K(B,2))$ (\code{H3Proof.rt-ptmap}, by induction on the two
pointed equivalences, \code{PtEquivOps.ptEqInd}). By
Corollary~\ref{cor:pointed}, $\Map(K(\Z,1),K(B,2))\simeq\Om K(B,2)\simeq
K(B,1)$ (\code{KK.e 1}), and $K(B,1)$ is connected (\code{KK.conn 0}).
So $\Map(\Om A_2,\Om B_3)$ is connected (\code{conn-rt}), and $\lambda f$, $\lambda 0$ are
merely equal.

\emph{Conclusion.} Given $q:\lambda f=\lambda 0$, the elements
$(f,\mathsf{refl})$ and $(0,q)$ of the proposition
$\code{Unloop}(\lambda f)$ are equal; their first components give
$f=0$. Since the goal is a proposition, the mere existence of $q$ suffices.
\end{proof}

\begin{remark}
(1) Only (K1), (K2) are used: the proof applies to any pointed connected
type with a homomorphic band onto $\Z$. (2) Theorem~\ref{thm:circle} is
stated for $Y$ in the universe of $K$; no other size condition appears.
(3) For (H4) the same looping has set-level fibres ($t=2$) and target
$\Map(K(\Z,1),K(B,3))\simeq\Om K(B,3)\simeq K(B,2)$; this is not treated
here.
\end{remark}

\section{Correspondence with the formal proof}\label{sec:corr}

\noindent Files \code{sag/pi3-s2-postnikov/CircleProperty.seg},
\code{H3Proof.seg}; \code{--opens} specification \code{circle.opens}.

\smallskip
\begin{tabular}{ll}
Lemma~\ref{lem:N} & \code{NatRec[W,psi,w].contr}\\
Lemma~\ref{lem:Z} & \code{IntRec[W,phi,w].contr} (with \code{equiv-rt}, \code{rt-data}, \code{res-equiv})\\
$(\ast)$ & \code{Circle.band-g} (then \code{g} is locked)\\
Lemma~\ref{lem:Q} & \code{Circle[Y].Pt[y,l].Q-contr} (\code{Q0-rt}, \code{cond-rt}, \code{rt-base'}, \code{pi-rt})\\
Theorem~\ref{thm:circle} (1)--(4) & \code{Pt.fib-rt} (\code{lp-rt}), \code{Pt.X1-rt} (\code{yon-rt}, \code{j-equiv}), \code{Pt.X2-rt}\\
Theorem~\ref{thm:circle} & \code{Circle[Y].E-equiv}\\
Corollary~\ref{cor:pointed} & \code{Circle[Y].evPt}, \code{evPt-conj}, \code{pointed-up}, \code{pointed-equiv}, \code{ptmap-conn}\\
 & (\code{fib-up}, \code{omega-up}, \code{subst-conj}, \code{sub-equiv})\\
Theorem~\ref{thm:H3} & \code{H3Proof.h3 : KZ2Cohomology.H3}
\end{tabular}

\end{document}
